% beamer2html.tex -- the manual of beamer2html (compile: pdflatex, twice)
%
% Copyright (C) 2026 Arun Periyal
% This work may be distributed and/or modified under the conditions of the
% LaTeX Project Public License, either version 1.3c of this license or (at
% your option) any later version (https://www.latex-project.org/lppl.txt).
\documentclass[a4paper,11pt]{article}
\usepackage[T1]{fontenc}
\usepackage{lmodern}
\usepackage{microtype}
\DisableLigatures[-]{family=tt*}   % --pack, not an en dash
\usepackage[margin=2.5cm]{geometry}
\usepackage{booktabs,tabularx,array}
\usepackage{xcolor}
\usepackage{fancyvrb}
\usepackage[hidelinks]{hyperref}

\newcommand\pkg{\textsf{beamer2html}}
\newcommand\cs[1]{\texttt{\textbackslash #1}}
\newcommand\env[1]{\texttt{#1}}
\newcommand\opt[1]{\texttt{#1}}
\newcommand\file[1]{\texttt{#1}}
\newcommand\key[1]{\fbox{\footnotesize\strut #1}}
\fvset{frame=single, framerule=0.4pt, fontsize=\small, xleftmargin=1em, xrightmargin=1em}
\newcolumntype{L}{>{\raggedright\arraybackslash}X}

\title{\pkg: beamer decks as offline HTML slide decks\\[0.5ex]
  \large version 1.0}
\author{Arun Periyal\\ \url{https://github.com/arunperiyal/tex2html}}
\date{2026-11-12}

\begin{document}
\maketitle

\begin{abstract}
\noindent \pkg{} turns a beamer presentation into a self-contained HTML slide
deck that runs offline in any modern browser, so animations, step-by-step
reveals, interactive plots, videos and speaker notes work without depending
on a particular PDF viewer. The same sources still give the PDF. It has two
parts: the LaTeX package \file{beamer2html.sty} (title slides, title bar,
footer, layout, tables, plots, animations, shapes, videos, flowcharts), used
in the beamer source, and the program \texttt{beamer2html} (Python), which
writes \file{slides.html} from that source.
\end{abstract}

\tableofcontents

\section{Introduction}

A beamer deck is usually presented from its PDF. That works until the deck
needs more than a PDF viewer offers everywhere: an animation (the
\textsf{animate} package needs Adobe Reader or a recent Okular), a video, a
plot whose values one wants to read off, or a presenter view with notes.
\pkg{} keeps the beamer source as the one source of the talk and adds an HTML
version of it:
\begin{itemize}
\item \file{slides.html} with \file{html/} next to it: a 16:9 deck that runs in
  any modern browser, offline, from a local file;
\item step-by-step reveals (\cs{pause}, overlay specifications) as in the PDF;
\item animations from image sequences, with play/pause, a time bar and
  synchronised panels;
\item interactive plots from data files (values under the pointer, zoom,
  hide lines), drawn by \textsf{pgfplots} in the PDF;
\item videos that play when their slide comes up;
\item a presenter view (notes, next slide, timer, slide grid; touch
  controls), a slide grid and fullscreen;
\item one zip with exactly the files the deck uses, for the presenting
  computer (\opt{--pack}), optionally with the images made smaller.
\end{itemize}

The program understands the LaTeX of decks written with
\file{beamer2html.sty}: its own commands and environments, beamer frames and
overlays, and common LaTeX (lists, figures, tables, math, citations). It does
not understand LaTeX in general. Anything it does not know is kept as plain
text and reported as a warning with the slide it is on, so nothing
disappears silently.

\section{Requirements and installation}

\begin{itemize}
\item Python 3.8 or later. Nothing else is needed for the conversion itself.
\item Poppler (\texttt{pdftocairo}, \texttt{pdftoppm}, \texttt{pdfimages},
  \texttt{pdfinfo}) for PDF figures: browsers cannot show PDF images, so
  vector figures become SVG and figures with embedded bitmaps PNG.
\item A TeX distribution (\texttt{pdflatex}, \texttt{pdfcrop}) and
  Ghostscript (\texttt{gs}) for TikZ pictures, which are drawn by LaTeX
  itself.
\item Optional: Pillow (\texttt{pip install pillow}) for
  \opt{--pack --compress}.
\item For the PDF: beamer, \textsf{tikz}, \textsf{pgfplots},
  \textsf{animate}, \textsf{biblatex} with \texttt{biber}, as usual.
\end{itemize}

\paragraph{From a TeX distribution.} When \pkg{} is installed by TeX Live or
MiKTeX, \file{beamer2html.sty} is found by LaTeX, and the program is the
command \texttt{beamer2html}.

\paragraph{From the CTAN package or the source repository.} Unpack it
anywhere. The program is the folder \file{scripts/beamer2html/} (in the
repository: its top folder); run it with Python:
\begin{Verbatim}
python3 path/to/beamer2html [main.tex] [options]
\end{Verbatim}
or through the launcher \file{beamer2html.py} in that folder, which can be
linked into a directory on the \texttt{PATH}. The program finds
\file{beamer2html.sty} next to itself (\file{latex/} in the repository,
\file{tex/} in the CTAN package) or through \texttt{kpsewhich}, and copies
it into each project it converts, so the PDF build finds it too; nothing has
to be installed into the TeX tree.

\section{Quick start}

The folder \file{example/} (in the repository: \file{template/}) is a
working deck that shows every element, one section each, with the options
explained in comments. Copy it to start a talk:
\begin{Verbatim}
cp -r path/to/example my-talk
cd my-talk
beamer2html              # -> slides.html (installs beamer2html.sty and html/)
latexmk -pdf main.tex    # -> main.pdf
\end{Verbatim}
Open \file{slides.html} in a browser. Fill in the \texttt{[placeholders]} in
\file{main.tex}, write your sections (copy the frames you need from the
tutorial ones), change the \cs{input} lines, and delete the tutorial
sections and the sample figures.

\section{The command line}

\begin{Verbatim}
beamer2html [main.tex] [-o slides.html] [--watch]
            [--pack [talk.zip] [--compress [PX]] [--slide-range RANGES]]
\end{Verbatim}
Run it in the folder of the main \file{.tex} file (default
\file{main.tex}); all paths are relative to that folder.

\medskip\noindent
\begin{tabularx}{\linewidth}{@{}lL@{}}
\toprule
Option & Meaning \\
\midrule
\opt{-o} \textit{file} & the HTML file to write (default \file{slides.html}) \\
\opt{-w}, \opt{--watch} & keep running and rebuild whenever \file{main.tex},
  an \cs{input} file, the \file{.bib}, an image or data file the slides use,
  or the tool's own files change. An open deck built this way reloads itself
  and stays on the same slide and step; a failed build leaves the last good
  deck in place. Ctrl+C stops it. Run \texttt{latexmk -pdf -pvc main.tex}
  alongside for the PDF. \\
\opt{--pack} [\textit{zip}] & also write a zip with the deck and exactly the
  images, animation frames, videos and \file{html/} files it uses (default
  name: the HTML file's name with \file{.zip}). It unpacks to one folder;
  open \file{slides.html} in it. Sources, the \file{.bib}, PDFs and unused
  figures are left out. \\
\opt{--compress} [\textit{PX}] & with \opt{--pack}: raster images and
  animation frames (PNG, JPEG, TIFF, BMP) go into the zip as WebP, at most
  \textit{PX} pixels wide and high (default 960). This often makes the zip
  5--10 times smaller with no visible difference; your files are not
  changed. Copies are cached in \file{html/cache/webp-PX/}. Needs Pillow. \\
\opt{--slide-range} \textit{R} & with \opt{--pack}: only these slides, e.g.\
  \texttt{20:30,32:35}; \texttt{5} is one slide, \texttt{20:} from 20 to the
  end, \texttt{:10} up to 10. The numbers are those the HTML deck shows; the
  packed deck numbers its slides from 1. \\
\opt{--version} & print the version \\
\bottomrule
\end{tabularx}

\medskip
Each run also installs into the project: \file{beamer2html.sty} next to the
main file; \file{html/deck.js}, \file{html/deck.css} and MathJax in
\file{html/}; \file{html/references.js} from the \cs{addbibresource} file;
and converted PDF figures and TikZ pictures in \file{html/cache/}. A copy of
\file{beamer2html.sty} that was edited in the project is saved as
\file{beamer2html.sty.bak} before it is replaced.

\section{Presenting}

\begin{tabularx}{\linewidth}{@{}lL@{}}
\toprule
Key & Action \\
\midrule
$\rightarrow$ $\downarrow$ Space PgDn & next step or slide \\
$\leftarrow$ $\uparrow$ PgUp & previous step or slide \\
Home, End; \textit{number} Enter & first, last slide; go to a slide \\
\key{F} & fullscreen \\
\key{P} & play / pause the animations and videos on the slide \\
\key{G} & slide grid: every slide as a thumbnail; click one to go there \\
\key{S} & presenter view in a second window \\
\key{H} & help \\
swipe, tap & on a touch screen: swipe left / right, or tap the right two
  thirds / left third \\
\bottomrule
\end{tabularx}

\medskip
The presenter view shows the current slide, the next step or slide, the
notes (\cs{note}), a clock and a timer. Put it on the laptop screen and the
slides on the projector; either window navigates. On a touch screen it has
buttons for back, play/pause, next and a slide grid that opens in the
presenter window only, so a slide can be picked while the audience keeps
seeing the current one; a swipe or a tap on the next-slide preview also
moves on. Printing (Ctrl+P, Save as PDF) gives a static copy, one page per
slide.

On a phone or tablet, browsers do not load the files next to a local HTML
file. The companion viewer (\url{https://arunperiyal.github.io/tex2html-viewer/},
also an Android app) opens a \opt{--pack} zip and presents it full screen,
offline.

\section{The \texttt{beamer2html} package}

Load it after beamer, together with the packages it works with:
\begin{Verbatim}
\documentclass[aspectratio=169,11pt]{beamer}
\usepackage{beamer2html}
\usepackage{animate}         % for the animation environment
\usepackage[backend=biber]{biblatex}  % for \ReferencesFrame and citations
\end{Verbatim}
The class font size option (8pt to 20pt) scales all text, in the PDF and
the HTML. The example deck shows every command below with its options in
comments; this section is an overview.

\subsection{Title slides and fields}

Set the fields in the preamble, then use a title slide:
\begin{Verbatim}
\SetTitle{...}  \SetSeminar{...}  \SetDegree{...}  \SetPresenter{...}
\SetRollNo{...} \SetSupervisor{...} \SetDepartment{...} \SetInstitute{...}
\SetLogo{figures/logo}          % without extension; a placeholder until it exists
\SetAuthorOne{...} \SetAuthorTwo{...} \SetConference{...}
\SetShortTitle{...} \SetFootAuthor{...}   % footer of content slides
\end{Verbatim}
\cs{SeminarTitleFrame} (thesis or seminar talk), \cs{ConferenceTitleFrame}
(two authors), \cs{OverviewFrame} (the \cs{section} names),
\cs{ReferencesFrame} (the bibliography; a long list continues on
``References (cont.)'' slides) and \cs{ThankYouFrame} make whole slides.

\subsection{Content slides}

Ordinary \verb|\begin{frame}{Title}{Subtitle}| frames get a title bar and a
footer. These switches work in the preamble and between frames:
\cs{TitleBarOn}/\cs{TitleBarOff}, \cs{SetTitleHeight}\verb|{len}|,
\cs{SectionNameOn}/\cs{SectionNameOff}, \cs{FooterOn}, \cs{FooterNumberOnly},
\cs{FooterOff}, \cs{SetTextMargins}\verb|{left}{right}| (the slide body;
beamer's \verb|\setbeamersize{text margin left=...}| is read too),
\cs{SetBarMargin}\verb|{len}| (title bar and footer),
\cs{SetCaptionSize}\verb|{\footnotesize}|, and
\cs{showdividers}/\cs{hidedividers}, an editing aid that outlines rows and
columns.

\subsection{Layout and tables}

The \env{rows} environment with \verb|\begin{row}{0.35}| (a share of the
slide height) holds beamer \env{columns}, giving a grid.
\env{slidetable} is a table in one of four styles, \opt{plain},
\opt{lines}, \opt{header} and \opt{striped}; its first row is the header,
and \texttt{p\{3cm\}}, \texttt{m\{..\}}, \texttt{b\{..\}} column widths carry
over to the HTML (\cs{slidetablestretch} sets the row spacing).
\env{progress} with \cs{task} lines is a timeline, \env{publications} with
\cs{pub}, \cs{pubtext} and \cs{pubsection} a list with status badges, and
\env{litreview} with \cs{study}, \cs{studytext} a literature table.
\cs{FigurePlaceholder} keeps the space of a figure that is not made yet.

\subsection{Shapes}

\verb|\begin{annotate}[grid] ... \end{annotate}| around an image or table,
with \verb|\shape<spec>[options]{kind}{x1,y1}{x2,y2}| lines, draws over it;
coordinates are fractions of the element, from the bottom left.
\cs{slideshape} uses fractions of the slide, and \cs{slidegrid} shows a grid
for reading them off. Kinds: \opt{rect}, \opt{ellipse}/\opt{circle},
\opt{line}, \opt{arrow}, \opt{text} (\verb|{x,y}{text}|) and \opt{axes} (a
coordinate indicator). Options: \opt{color}, \opt{width}, \opt{fill},
\opt{opacity}, \opt{dashed}, \opt{rounded}, \opt{label}, \opt{labelpos},
\opt{textsize}, \opt{border}. Shapes never change the layout.

\subsection{Videos}

\verb|\slidevideo[options]{file}| (MP4, WebM) with \opt{width},
\opt{height}, \opt{poster}, \opt{autoplay} (when its slide or step comes
up), \opt{loop}, \opt{muted}, \opt{controls}, \opt{start} and \opt{end} in
seconds; \verb|\slidevideo<2->| for a step. In the PDF it is the poster (or
a box with the file name) linking to the video, or with
\verb|\usepackage{multimedia}| a \cs{movie}.

\subsection{Plots}

\begin{Verbatim}
\begin{slideplot}[xlabel={$t$}, ylabel={$y/D$}, height=5cm]
  \plotline[x=time, y=disp, label={$U_r = 4$}]{data/U4.csv}
  \plotline<2->[x=time, y=disp, label={$U_r = 6$}, style=dashed]{data/U6.csv}
\end{slideplot}
\end{Verbatim}
draws lines from data files (CSV or space-separated columns, optional header
row; columns by name or by number from 1). The PDF uses \textsf{pgfplots};
the HTML plot is interactive. Options of \env{slideplot}: \opt{width},
\opt{height}, \opt{xlabel}, \opt{ylabel}, \opt{xmin}, \opt{xmax},
\opt{ymin}, \opt{ymax}, \opt{grid}, \opt{legend} (\opt{north east},
\opt{north west}, \opt{south east}, \opt{south west}, \opt{outer east},
\opt{outer north}, \opt{none}), \opt{rows}, \opt{columns}, \opt{share}. Per
line: \opt{x}, \opt{y}, \opt{label}, \opt{color}, \opt{style}
(\opt{solid}, \opt{dashed}, \opt{dotted}), \opt{xscale}, \opt{yscale}
(multiply the values), \opt{offset} (added to x, after \opt{xscale}),
\opt{sep}, \opt{pdfevery=N} (the PDF plots every $N$-th point). Lines with
the same label share one colour and one legend entry.

\verb|\plotpanel[title=..., xlabel=..., ylabel=...]| starts a panel; the
lines after it are drawn in it. \opt{rows}/\opt{columns} lay the panels
out, \opt{share} (\opt{x}, \opt{y}, \opt{both}, \opt{none}) links axes, and
a panel's own \opt{xmin}, \opt{xmax}, \opt{ymin}, \opt{ymax} apply to an
axis it does not share. There is one legend for all panels.

\subsection{Animations}

\begin{Verbatim}
\begin{animation}[frames=0:22:2, fps=3, controls=hover, dt=0.1, timebar]
  \animpanel{Case 1}{figures/anim/case1_\N}
  \animpanel{Case 2}{figures/anim/case2_\N}
\end{animation}
\end{Verbatim}
plays image sequences in synchronised panels (\cs{N} is the frame number).
Options: \opt{frames=first:last:step}, \opt{fps}, \opt{width},
\opt{controls} (including \opt{hover}), \opt{autoplay}, \opt{loop},
\opt{subcaptions}, \opt{dt} with \opt{timelabel}, \opt{timedigits},
\opt{timeunit}, \opt{time=false}, \opt{timebar}, \opt{columns=N}; a panel
can have its own \verb|\animpanel[frames=a:b:s, dt=...]|. In the PDF it is
an \textsf{animate} animation. \textsf{animate}'s own \env{animateinline}
with \cs{multiframe} is understood too.

\subsection{Diagrams}

\env{tikzpicture} and \cs{tikz} are drawn by LaTeX itself: each picture is
compiled on its own with the deck's preamble, cropped and turned into one
SVG per overlay step, so the HTML matches the PDF; overlays (\verb|\node<2->|,
\verb|visible on=<2->|, \cs{pause} inside the picture) work as in beamer.
The package has flowchart styles in the deck's colours: \opt{flow} on the
picture, nodes \opt{flowbox}, \opt{flowstart}, \opt{flowdecision},
\opt{flowio}, and \opt{flowarrow}/\opt{flowlabel}.

\section{What the conversion understands}

\begin{itemize}
\item beamer: frames (\opt{[t]}, \opt{[plain]}, \cs{frametitle}),
  \cs{section}, \cs{pause}, overlay specifications (\verb|\item<2->|,
  \verb|[<+->]|, \cs{only}, \cs{uncover}, \cs{visible}, \cs{invisible},
  \cs{alt}, \cs{onslide}, environments with \verb|<...>|), \env{columns},
  \env{block}, \cs{note}, the class font size option;
\item LaTeX: \env{figure} and \env{table} with \cs{caption} (numbered through
  the deck), \cs{label} and \cs{ref}, lists, \env{minipage}, \env{center},
  \env{tabular}, math (MathJax~4, in a Times font when the preamble loads
  \textsf{newtxmath}, \textsf{mathptmx}, \textsf{txfonts} or \textsf{times},
  else New Computer Modern), \cs{includegraphics}, font sizes and styles,
  \cs{color}, \cs{cite}, \cs{textcite}, \cs{citet}, \cs{citeauthor},
  \cs{href}, \cs{def} and \cs{newcommand} without arguments;
\item everything in \file{beamer2html.sty} described above.
\end{itemize}
Not supported: commands with arguments defined in the document,
\verb|\alert<...>|.

\section{Decks from before the rename}

\pkg{} was called \textsf{tex2html}, and its package \textsf{titletemplates}.
A deck that loads \verb|\usepackage{titletemplates}| keeps working: its
\file{titletemplates.sty} is updated together with \file{beamer2html.sty}.
To move it over, change the line to \verb|\usepackage{beamer2html}|.

\section{License}

Copyright \copyright{} 2026 Arun Periyal. This work (the program,
\file{deck.js}, \file{deck.css}, \file{beamer2html.sty}, the example deck
and this manual) may be distributed and/or modified under the conditions of
the \LaTeX{} Project Public License, either version 1.3c of this license or
(at your option) any later version. It has the LPPL maintenance status
``maintained''; the Current Maintainer is Arun Periyal.

Bundled third-party files keep their own licenses: MathJax and its fonts
(Apache License 2.0) and a subset of TeX Gyre Termes as WOFF2 (GUST Font
License), both in the program's \file{assets/vendor/} folder with their
license files.

\section{Changes}

\begin{description}
\item[1.0 (2026-11-12)] First release on CTAN, as \pkg{} (formerly
  \textsf{tex2html} with \textsf{titletemplates.sty}).
\end{description}

\end{document}
